\documentclass[11pt]{amsart}
\usepackage{geometry}                % See geometry.pdf to learn the layout options. There are lots.
\geometry{letterpaper}                   % ... or a4paper or a5paper or ... 
%\geometry{landscape}                % Activate for for rotated page geometry
%\usepackage[parfill]{parskip}    % Activate to begin paragraphs with an empty line rather than an indent
\usepackage{graphicx}
\usepackage{amssymb}
\usepackage{epstopdf}
\usepackage{multirow}
\DeclareGraphicsRule{.tif}{png}{.png}{`convert #1 `dirname #1`/`basename #1 .tif`.png}

\title{ Ejemplo 2.1}
\author{Enrique Zuñiga Zuñiga}
%\date{}                                           % Activate to display a given date or no date

\begin{document}
\maketitle
%\section{}
Ejemplo 2.1\newline
Construya un diagrama de flujo tal que dado como dato la calificación de un alumno en un examen, escriba “aprobado” en caso de que esa calificación sea mayor a 8.\newline\newline
Dato: CAL (variable de tipo real que representa la calificación del alumno).\newline\newline
En la tabla 2.1, podemos observar el seguimiento del diagrama de flujo para
diferentes corridas.\newline


$$
\begin{array}{|c|c|c|}
\hline
 \multicolumn{3}{|c|}{Tabla 2.1}\\
\hline
NUMERO DECORRIDA & DATO & RESULTADO \\
 & CAL &  \\
\hline
1 & 8.75 & \text{``Aprobado''} \\
\hline
2 & 7.90 &  \\
\hline
3 & 8.00 &  \\
\hline
4 & 9.50 & \text{``Aprobado''} \\
\hline
5 & 8.35 & \text{``Aprobado''} \\
\hline
\end{array}
$$\newline\newline
A continuación presentamos el diagrama de flujo en lenguaje algorítmico.\newline\newline
EXAMEM SELECTIVA SIMPLE\newline
El programa, dado com o dato la calificación de un alumno en un examen, escribe aprobado si la calificación es superior a 8 CAL es una variable de tipo real\newline
1. Leer CAL\newline
2. S1 CAL > 8 entonces\newline
Escribir "Aprobado"\newline
3. Fin del condicional d e l paso 2 \newline

%\subsection{}



\end{document}  